% GENERATED FILE. Edit canonical lessons, then run npm run generate:books.
% canonical-source-hash: fnv1a64:4e245a6ef74e66ce
% canonical-lessons: ML-C291-valli, ML-C291-jantu, ML-C291-man, ML-C291-kala3, ML-C291-cennaya

\chapter{Creeper, Animal, Deer, Bull, Wolf}
\label{ch:ml-creeper-animal-deer-bull-wolf}
\hlchaptermodality{291}

\begin{chapteropening}
\textbf{By the end of this chapter:} \emph{I can say a creeper, an animal, a deer, a bull and a wolf.}

Complete the last lesson of chapter 291: I can say a creeper, an animal, a deer, a bull and a wolf.
\end{chapteropening}

\section[\texorpdfstring{\ml{വള്ളി} (vaḷḷi)}{വള്ളി (vaḷḷi)}]{\ml{വള്ളി} (vaḷḷi) — a creeper, a vine}

\label{lesson:ML-C291-valli}



\begin{quote}
Before the new one: say the Malayalam for a chenda drum, then the Malayalam for a pooram.
\end{quote}

\subsection*{You'll want to know: \ml{വള്ളി}}
\textbf{\ml{വള്ളി}} — \emph{vaḷḷi} — \textquotedblleft{}a creeper, a vine\textquotedblright{}.

\begin{grammarlens}[title={What you've built}]
One more word for this chapter.
\end{grammarlens}

\subsection*{Your turn}
\begin{itemize}
\raggedright
  \item \emph{Say it:} \emph{vaḷḷi}
  \item \emph{Say it:} \emph{vaḷḷi}, once more
  \item \emph{Say it:} say \emph{vaḷḷi}
  \item \emph{Recall:} say the Malayalam for a chenda drum, then the Malayalam for a pooram, then say \emph{vaḷḷi} again
\end{itemize}

\begin{tcolorbox}[breakable,colback=teal!4,colframe=teal!35!black,title={Before you move on}]
What does \ml{വള്ളി} mean? (\textquotedblleft{}A creeper, a vine\textquotedblright{}.) Say it once more.
\end{tcolorbox}
\section[\texorpdfstring{\ml{ജന്തു} (jantu)}{ജന്തു (jantu)}]{\ml{ജന്തു} (jantu) — an animal}

\label{lesson:ML-C291-jantu}



\begin{quote}
Before the new one: say the Malayalam for a pooram, then the Malayalam for a creeper.
\end{quote}

\subsection*{You'll want to know: \ml{ജന്തു}}
\textbf{\ml{ജന്തു}} — \emph{jantu} — \textquotedblleft{}an animal\textquotedblright{}.

\begin{grammarlens}[title={What you've built}]
One more word for this chapter.
\end{grammarlens}

\subsection*{Your turn}
\begin{itemize}
\raggedright
  \item \emph{Say it:} \emph{jantu}
  \item \emph{Say it:} \emph{jantu}, once more
  \item \emph{Say it:} say \emph{jantu}
  \item \emph{Recall:} say the Malayalam for a pooram, then the Malayalam for a creeper, then say \emph{jantu} again
\end{itemize}

\begin{tcolorbox}[breakable,colback=teal!4,colframe=teal!35!black,title={Before you move on}]
What does \ml{ജന്തു} mean? (\textquotedblleft{}An animal\textquotedblright{}.) Say it once more.
\end{tcolorbox}
\section[\texorpdfstring{\ml{മാൻ} (mān)}{മാൻ (mān)}]{\ml{മാൻ} (mān) — a deer}

\label{lesson:ML-C291-man}



\begin{quote}
Before the new one: say the Malayalam for a creeper, then the Malayalam for an animal.
\end{quote}

\subsection*{You'll want to know: \ml{മാൻ}}
\textbf{\ml{മാൻ}} — \emph{mān} — \textquotedblleft{}a deer\textquotedblright{}.

\begin{grammarlens}[title={What you've built}]
One more word for this chapter.
\end{grammarlens}

\subsection*{Your turn}
\begin{itemize}
\raggedright
  \item \emph{Say it:} \emph{mān}
  \item \emph{Say it:} \emph{mān}, once more
  \item \emph{Say it:} say \emph{mān}
  \item \emph{Recall:} say the Malayalam for a creeper, then the Malayalam for an animal, then say \emph{mān} again
\end{itemize}

\begin{tcolorbox}[breakable,colback=teal!4,colframe=teal!35!black,title={Before you move on}]
What does \ml{മാൻ} mean? (\textquotedblleft{}A deer\textquotedblright{}.) Say it once more.
\end{tcolorbox}
\section[\texorpdfstring{\ml{കാള} (kāḷa)}{കാള (kāḷa)}]{\ml{കാള} (kāḷa) — a bull, an ox}

\label{lesson:ML-C291-kala3}



\begin{quote}
Before the new one: say the Malayalam for an animal, then the Malayalam for a deer.
\end{quote}

\subsection*{You'll want to know: \ml{കാള}}
\textbf{\ml{കാള}} — \emph{kāḷa} — \textquotedblleft{}a bull, an ox\textquotedblright{}.

\begin{grammarlens}[title={What you've built}]
One more word for this chapter.
\end{grammarlens}

\subsection*{Your turn}
\begin{itemize}
\raggedright
  \item \emph{Say it:} \emph{kāḷa}
  \item \emph{Say it:} \emph{kāḷa}, once more
  \item \emph{Say it:} say \emph{kāḷa}
  \item \emph{Recall:} say the Malayalam for an animal, then the Malayalam for a deer, then say \emph{kāḷa} again
\end{itemize}

\begin{tcolorbox}[breakable,colback=teal!4,colframe=teal!35!black,title={Before you move on}]
What does \ml{കാള} mean? (\textquotedblleft{}A bull, an ox\textquotedblright{}.) Say it once more.
\end{tcolorbox}
\section[\texorpdfstring{\ml{ചെന്നായ} (cennāya)}{ചെന്നായ (cennāya)}]{\ml{ചെന്നായ} (cennāya) — a wolf}

\label{lesson:ML-C291-cennaya}



\begin{quote}
Before the new one: say the Malayalam for a deer, then the Malayalam for a bull.
\end{quote}

\subsection*{You'll want to know: \ml{ചെന്നായ}}
\textbf{\ml{ചെന്നായ}} — \emph{cennāya} — \textquotedblleft{}a wolf\textquotedblright{}.

\begin{grammarlens}[title={What you've built}]
One more word for this chapter.
\end{grammarlens}

\subsection*{Your turn}
\begin{itemize}
\raggedright
  \item \emph{Say it:} \emph{cennāya}
  \item \emph{Say it:} \emph{cennāya}, once more
  \item \emph{Say it:} say \emph{cennāya}
  \item \emph{Recall:} say the Malayalam for a deer, then the Malayalam for a bull, then say \emph{cennāya} again
\end{itemize}

\begin{tcolorbox}[breakable,colback=teal!4,colframe=teal!35!black,title={Before you move on}]
What does \ml{ചെന്നായ} mean? (\textquotedblleft{}A wolf\textquotedblright{}.) Say it once more.
\end{tcolorbox}
